% test-014.tex -- comandos \spangle, \spmangle, \spang, \sparc y \spmarc
%
% Compilar:  lualatex-dev test-014.tex
% Validar:   show-pdf-tags --xml test-014.pdf | rnv-wrapp
%            verapdf --flavour ua2 --format text test-014.pdf
%
% Cada caso es un parrafo numerado, para poder referirse a el al escuchar
% el PDF. Los comentarios "doc:" indican lo que documenta el manual; en el
% resto, la lectura se revisa escuchando. Todos trabajan en modo
% matematico. \spangle y \spmangle reciben un punto, una letra griega, un
% numero o tres puntos; \sparc y \spmarc, dos o tres puntos; \spang, un
% angulo sexagesimal y no tiene llaves. Todos los casos de \spang
% requieren spang-expand.patch, y los de solo grados (45, 180 y el 90
% de la formula) ademas spang-grados.patch. Las entradas invalidas
% (que detienen la compilacion) no van aqui. \spangle acepta sin error
% entradas que el manual no documenta (AB, ABCD, a, x, 1.5): tampoco van.
%
% Casos pensados para escuchar: prefix con una letra griega o un numero en
% \spangle y \spmangle, y concept=false en \sparc y \spmarc
% (sparc-concept.patch).
\DocumentMetadata
  {
    lang=es-CL,
    pdfversion=2.0,
    pdfstandard=ua-2,
    tagging=on,
    tagging-setup={math/setup={mathml-SE}},
  }
\documentclass{article}
\usepackage[spanish]{babel}
\usepackage{unicode-math}
\usepackage{spintent}
\usepackage{hyperref}
\hypersetup
  {
    pdfdisplaydoctitle = true,
    pdftitle           = {test-014: spangle, spmangle, spang, sparc y spmarc},
  }
\begin{document}

\section*{Comando spangle: argumentos}

% doc: ángulo en A
1. Un punto: $\spangle{A}$.

% doc: ángulo alfa
2. Una letra griega: $\spangle{\alpha}$.

% doc: ángulo uno
3. Un número: $\spangle{1}$.

% doc: ángulo A O B
4. Tres puntos: $\spangle{AOB}$.

5. Un punto con prima: $\spangle{A'}$.

6. Un punto con subíndice: $\spangle{A_{1}}$.

7. Otro punto: $\spangle{B}$.

8. Con la letra beta: $\spangle{\beta}$.

9. Con la letra gamma: $\spangle{\gamma}$.

10. Con la letra theta: $\spangle{\theta}$.

11. Con la letra varphi: $\spangle{\varphi}$.

12. Un número de dos cifras: $\spangle{12}$.

13. Tres puntos con comas: $\spangle{A,O,B}$.

14. Tres puntos con comas y espacios: $\spangle{A, O, B}$.

15. Con primas: $\spangle{A'OB'}$.

16. Con subíndices: $\spangle{A_{1}OB_{2}}$.

17. Con subíndices y comas: $\spangle{A_{1},O,B_{2}}$.

18. Con otras letras: $\spangle{PQR}$.

19. Con tres vértices: $\spangle{ABC}$.

\section*{Comando spangle: llaves}

% doc: sin el prefijo mayúscula
20. Con read-cmd school, el valor por defecto: $\spangle[read-cmd=school]{AOB}$.

% doc: MathCAT aplica sus propias reglas
21. Con read-cmd mathcat: $\spangle[read-cmd=mathcat]{AOB}$.

22. Con read-cmd mathcat y un punto: $\spangle[read-cmd=mathcat]{A}$.

% doc: el prefijo se verbaliza antes del concepto
23. Con prefix y un punto: $\spangle[prefix={el}]{A}$.

% doc: el prefijo se verbaliza antes del concepto
24. Con prefix y tres puntos: $\spangle[prefix={el}]{AOB}$.

% doc: el prefijo se verbaliza antes del concepto
25. Con prefix y una letra griega: $\spangle[prefix={el}]{\alpha}$.

% doc: el prefijo se verbaliza antes del concepto
26. Con prefix y un número: $\spangle[prefix={el}]{1}$.

% doc: ángulo recto en A
27. Con recto true y un punto: $\spangle[recto=true]{A}$.

% doc: ángulo recto A O B
28. Con recto true y tres puntos: $\spangle[recto=true]{AOB}$.

29. Con recto true y una letra griega: $\spangle[recto=true]{\alpha}$.

30. Con recto true y un número: $\spangle[recto=true]{1}$.

31. Con recto false, el valor por defecto: $\spangle[recto=false]{AOB}$.

32. Con concept true, el valor por defecto: $\spangle[concept=true]{AOB}$.

% doc: solo se verbaliza el argumento
33. Con concept false y un punto: $\spangle[concept=false]{A}$.

% doc: solo se verbaliza el argumento
34. Con concept false y tres puntos: $\spangle[concept=false]{AOB}$.

% doc: solo se verbaliza el argumento
35. Con concept false y una letra griega: $\spangle[concept=false]{\alpha}$.

% doc: solo se verbaliza el argumento
36. Con concept false y un número: $\spangle[concept=false]{1}$.

37. Con silent false, el valor por defecto: $\spangle[silent=false]{AOB}$.

% doc: no se verbaliza
38. Con silent true y un punto: $\spangle[silent=true]{A}$.

% doc: no se verbaliza
39. Con silent true y tres puntos: $\spangle[silent=true]{AOB}$.

% doc: no se verbaliza
40. Con silent true y una letra griega: $\spangle[silent=true]{\alpha}$.

% doc: no se verbaliza
41. Con silent true y un número: $\spangle[silent=true]{1}$.

42. Con todas las llaves: $\spangle[read-cmd=mathcat,prefix={el},recto=true,concept=true,silent=false]{A_{1}OB_{2}}$.

\section*{Comando spmangle: argumentos}

% doc: medida del ángulo en A
43. Un punto: $\spmangle{A}$.

% doc: medida del ángulo alfa
44. Una letra griega: $\spmangle{\alpha}$.

% doc: medida del ángulo uno
45. Un número: $\spmangle{1}$.

% doc: medida del ángulo A O B
46. Tres puntos: $\spmangle{AOB}$.

47. Un punto con prima: $\spmangle{A'}$.

48. Un punto con subíndice: $\spmangle{A_{1}}$.

49. Otro punto: $\spmangle{B}$.

50. Con la letra beta: $\spmangle{\beta}$.

51. Con la letra gamma: $\spmangle{\gamma}$.

52. Con la letra theta: $\spmangle{\theta}$.

53. Con la letra varphi: $\spmangle{\varphi}$.

54. Un número de dos cifras: $\spmangle{12}$.

55. Tres puntos con comas: $\spmangle{A,O,B}$.

56. Tres puntos con comas y espacios: $\spmangle{A, O, B}$.

57. Con primas: $\spmangle{A'OB'}$.

58. Con subíndices: $\spmangle{A_{1}OB_{2}}$.

59. Con subíndices y comas: $\spmangle{A_{1},O,B_{2}}$.

60. Con otras letras: $\spmangle{PQR}$.

61. Con tres vértices: $\spmangle{ABC}$.

\section*{Comando spmangle: llaves}

% doc: sin el prefijo mayúscula
62. Con read-cmd school, el valor por defecto: $\spmangle[read-cmd=school]{AOB}$.

% doc: MathCAT aplica sus propias reglas
63. Con read-cmd mathcat: $\spmangle[read-cmd=mathcat]{AOB}$.

64. Con read-cmd mathcat y un punto: $\spmangle[read-cmd=mathcat]{A}$.

% doc: el prefijo se verbaliza antes del concepto
65. Con prefix y un punto: $\spmangle[prefix={el}]{A}$.

% doc: el prefijo se verbaliza antes del concepto
66. Con prefix y tres puntos: $\spmangle[prefix={el}]{AOB}$.

% doc: el prefijo se verbaliza antes del concepto
67. Con prefix y una letra griega: $\spmangle[prefix={el}]{\alpha}$.

% doc: el prefijo se verbaliza antes del concepto
68. Con prefix y un número: $\spmangle[prefix={el}]{1}$.

% doc: medida del ángulo recto en A
69. Con recto true y un punto: $\spmangle[recto=true]{A}$.

% doc: medida del ángulo recto A O B
70. Con recto true y tres puntos: $\spmangle[recto=true]{AOB}$.

71. Con recto true y una letra griega: $\spmangle[recto=true]{\alpha}$.

72. Con recto true y un número: $\spmangle[recto=true]{1}$.

73. Con recto false, el valor por defecto: $\spmangle[recto=false]{AOB}$.

74. Con concept true, el valor por defecto: $\spmangle[concept=true]{AOB}$.

% doc: solo se verbaliza el argumento
75. Con concept false y un punto: $\spmangle[concept=false]{A}$.

% doc: solo se verbaliza el argumento
76. Con concept false y tres puntos: $\spmangle[concept=false]{AOB}$.

% doc: solo se verbaliza el argumento
77. Con concept false y una letra griega: $\spmangle[concept=false]{\alpha}$.

% doc: solo se verbaliza el argumento
78. Con concept false y un número: $\spmangle[concept=false]{1}$.

79. Con silent false, el valor por defecto: $\spmangle[silent=false]{AOB}$.

% doc: no se verbaliza
80. Con silent true y un punto: $\spmangle[silent=true]{A}$.

% doc: no se verbaliza
81. Con silent true y tres puntos: $\spmangle[silent=true]{AOB}$.

% doc: no se verbaliza
82. Con silent true y una letra griega: $\spmangle[silent=true]{\alpha}$.

% doc: no se verbaliza
83. Con silent true y un número: $\spmangle[silent=true]{1}$.

% doc: imprime \measuredangle AOB
84. Con print-m false, el valor por defecto: $\spmangle[print-m=false]{AOB}$.

% doc: imprime m\angle AOB -- «medida del ángulo A O B»
85. Con print-m true: $\spmangle[print-m=true]{AOB}$.

86. Con print-m true y recto true: $\spmangle[print-m=true,recto=true]{AOB}$.

87. Con print-m true y un número: $\spmangle[print-m=true]{1}$.

88. Con todas las llaves: $\spmangle[read-cmd=mathcat,prefix={el},recto=true,concept=true,silent=false,print-m=true]{A_{1}OB_{2}}$.

\section*{Comando spang}

% doc: imprime 45°30′15″ -- «cuarenta y cinco grados treinta minutos quince segundos»
89. Con grados, minutos y segundos: $\spang{45:30:15}$.

90. Con grados y minutos: $\spang{45:30}$.

% doc: requiere spang-grados.patch
91. Solo grados: $\spang{45}$.

% doc: requiere spang-grados.patch
92. Solo grados, ángulo llano: $\spang{180}$.

93. Un segundo: $\spang{0:0:1}$.

94. Ángulo recto: $\spang{90:0:0}$.

95. Con grados de tres cifras y minutos: $\spang{120:15}$.

96. Una vuelta completa: $\spang{360:0:0}$.

97. Con unos y dos: $\spang{1:2:3}$.

98. Con cifras sueltas: $\spang{7:5}$.

99. Suma de ángulos: $\spang{45:30:15} + \spang{10:20:30}$.

% doc: requiere spang-grados.patch
100. Medida igual a un ángulo: $\spmangle{AOB} = \spang{90}$.

\section*{Comando sparc}

% doc: arco A B
101. Dos puntos juntos: $\sparc{AB}$.

102. Dos puntos con coma: $\sparc{A,B}$.

103. Dos puntos con coma y espacio: $\sparc{A, B}$.

104. Tres puntos: $\sparc{ABC}$.

105. Tres puntos con comas: $\sparc{A,B,C}$.

106. Con primas: $\sparc{A'B'}$.

107. Con subíndices: $\sparc{A_{1}B_{2}}$.

108. Tres puntos con subíndices: $\sparc{A_{1}B_{2}C_{3}}$.

109. Con el centro en medio: $\sparc{AOB}$.

110. Con read-cmd school, el valor por defecto: $\sparc[read-cmd=school]{AB}$.

% doc: MathCAT aplica sus propias reglas
111. Con read-cmd mathcat: $\sparc[read-cmd=mathcat]{AB}$.

% doc: el prefijo se verbaliza antes del concepto
112. Con prefix y dos puntos: $\sparc[prefix={el}]{AB}$.

113. Con prefix y tres puntos: $\sparc[prefix={el}]{ABC}$.

114. Con concept true, el valor por defecto: $\sparc[concept=true]{AB}$.

% doc: solo se verbaliza el argumento: A B
115. Con concept false y dos puntos: $\sparc[concept=false]{AB}$.

116. Con concept false y tres puntos: $\sparc[concept=false]{ABC}$.

117. Con concept false y subíndices: $\sparc[concept=false]{A_{1}B_{2}}$.

118. Con silent false, el valor por defecto: $\sparc[silent=false]{AB}$.

% doc: no se verbaliza
119. Con silent true y dos puntos: $\sparc[silent=true]{AB}$.

120. Con silent true y tres puntos: $\sparc[silent=true]{ABC}$.

121. Con silent true y subíndices: $\sparc[silent=true]{A_{1}B_{2}}$.

122. Con todas las llaves: $\sparc[read-cmd=mathcat,prefix={el},concept=true,silent=false]{A_{1}B_{2}C_{3}}$.

\section*{Comando spmarc}

% doc: medida del arco A B
123. Dos puntos juntos: $\spmarc{AB}$.

124. Dos puntos con coma: $\spmarc{A,B}$.

125. Dos puntos con coma y espacio: $\spmarc{A, B}$.

126. Tres puntos: $\spmarc{ABC}$.

127. Tres puntos con comas: $\spmarc{A,B,C}$.

128. Con primas: $\spmarc{A'B'}$.

129. Con subíndices: $\spmarc{A_{1}B_{2}}$.

130. Tres puntos con subíndices: $\spmarc{A_{1}B_{2}C_{3}}$.

131. Con el centro en medio: $\spmarc{AOB}$.

132. Con read-cmd school, el valor por defecto: $\spmarc[read-cmd=school]{AB}$.

% doc: MathCAT aplica sus propias reglas
133. Con read-cmd mathcat: $\spmarc[read-cmd=mathcat]{AB}$.

% doc: el prefijo se verbaliza antes del concepto
134. Con prefix y dos puntos: $\spmarc[prefix={el}]{AB}$.

135. Con prefix y tres puntos: $\spmarc[prefix={el}]{ABC}$.

136. Con concept true, el valor por defecto: $\spmarc[concept=true]{AB}$.

% doc: solo se verbaliza el argumento: A B
137. Con concept false y dos puntos: $\spmarc[concept=false]{AB}$.

138. Con concept false y tres puntos: $\spmarc[concept=false]{ABC}$.

139. Con concept false y subíndices: $\spmarc[concept=false]{A_{1}B_{2}}$.

140. Con silent false, el valor por defecto: $\spmarc[silent=false]{AB}$.

% doc: no se verbaliza
141. Con silent true y dos puntos: $\spmarc[silent=true]{AB}$.

142. Con silent true y tres puntos: $\spmarc[silent=true]{ABC}$.

143. Con silent true y subíndices: $\spmarc[silent=true]{A_{1}B_{2}}$.

144. Con todas las llaves: $\spmarc[read-cmd=mathcat,prefix={el},concept=true,silent=false]{A_{1}B_{2}C_{3}}$.

\section*{En fórmulas}

145. Ángulos iguales: $\spangle{ABC} = \spangle{DEF}$.

146. Suma de medidas: $\spmangle{AOB} + \spmangle{BOC} = \spmangle{AOC}$.

147. Medida de un ángulo en grados: $\spmangle{AOB} = \spang{45:30}$.

148. Arcos iguales: $\sparc{AB} = \sparc{CD}$.

149. Medida de un arco en grados: $\spmarc{AB} = \spang{45:30:15}$.

150. Ángulo y arco que lo subtiende: $\spmangle{AOB} = \spmarc{AB}$.

151. En fórmula destacada:
\[ \spmangle[print-m=true]{AOB} = \spmarc{AB} \].

\end{document}
